\section{실험 (Experiments)}
\label{sec:experiments}

본 절에서는 ScopeAD의 few-shot 성능을 비교하고, 관측 확장과 결합 설계가 결과에 미치는 영향을 분석한다. 주 비교에 이어 관측 구성과 계산 비용, 결합 대안, 획득 조건별 진단, 합성 열화 비교, 공식 private 평가 순으로 근거를 제시한다.

\subsection{평가 설정 (Evaluation Setup)}
\label{sec:experiments:setup}

\noindent\textbf{데이터셋 및 분할.}
MVTec AD 2~\cite{mvtecad2}는 실제 산업 환경의 조명 드리프트와 미세 결함을 포괄하는 8개 클래스 벤치마크다.
공개 개발 평가 세트인 \texttt{test\_public}은 총 1,084장(정상 379장, 결함 705장)으로 구성되며, 기준(Regular) 환경 외에 Shift 1--3 등의 획득 조건 변화를 포함한다.
정상 서포트 이미지는 \texttt{train/good}에서 고정 시드와 매니페스트로 선별한다.
교차 데이터셋 검증을 위해 기존 MVTec AD(15개 클래스)~\cite{mvtecad}와 VisA(12개 클래스)~\cite{visa} 벤치마크를 함께 평가한다.

\noindent\textbf{Few-shot 설정 및 비교군.}
클래스당 정상 서포트 이미지 수 $N$은 $\{1, 2, 4, 8, 16\}$으로 설정한다.
ScopeAD(Ours, linear0125)는 DINOv2 ViT-L/14 백본을 동결 사용하며, 계층 $\{5, 11, 23\}$과 4개 해상도 $\{448, 672, 896, 1120\}$, FP16 autocast, cap256, $z_d + 0.125\,z_5$ 결합을 적용한다.
비교 대상 방법군은 PatchCore~\cite{patchcore}, SuperAD~\cite{superad}, AnomalyDINO~\cite{anomalydino}, FastRecon~\cite{fastrecon}, RAD~\cite{rad}, SubspaceAD~\cite{subspacead}, RegAD~\cite{regad}다.
동일한 shot과 seed의 비교는 서포트 경로, 쿼리 및 평가기를 공유한다. 비교 파이프라인의 해상도 설정값은 1120을 기준으로 하며, ScopeAD는 짧은 변 길이의 최댓값을 1120으로 설정하고 종횡비를 유지한다. RegAD는 STN 제약으로 공식 224 설정을 유지한다. 이 비교는 백본, 전체 토큰 수 또는 계산량까지 일치시킨 비교는 아니다.
주요 평가 지표는 클래스 매크로 \AUPRO, \pAUROC, 그리고 이미지 단위 \IAUROC다.
각 클래스에서 AUPRO는 결함 연결영역별 검출 비율을 평균한 PRO 곡선을 FPR $[0,0.05]$에서 적분하고 0.05로 나눈다.
pAUROC는 같은 구간의 pixel ROC 면적에 무작위 기준 0.5, 완전 분류 기준 1의 표준화를 적용한다.
이미지 AUROC와 교차 데이터셋 표의 pixel AUROC는 전체 FPR 구간을 사용하며, 모든 지표는 클래스별 값을 동일 가중 평균한다.


\noindent\textbf{개발 조건과 실행 규약.}
결합 가중치는 AD2 \texttt{test\_public} 2-shot seed 0에서 선택했으므로, AD2 공개 결과는 개발에 사용된 벤치마크의 평가로 해석한다.
MVTec AD와 VisA에는 동일한 가중치를 적용한다. 각 평가 클래스의 뱅크 및 척도 추정에는 정상 서포트만 사용한다.
ScopeAD는 \cref{sec:method:aggregation}의 고점수 연결영역 후처리를 포함한다.
AD2 1/4/8-shot 반복 평가와 16-shot 용량 비교의 실행 기록에서, 잡음 진단에 따른 저해상도 전환은 모든 쿼리에서 비활성 상태였음을 확인했다.
따라서 이 실행에서는 각 구성에 지정된 해상도 집합을 모든 쿼리에 고정하여 사용했다.
% Provenance for refreshed metrics: results/remote_runs/dsba3/scopead_v18_q95_linear0125_fp16_20261012/combined/{strata_ad2,multiseed_ad2,capacity_cost_ad2}.
% Each refreshed macro was recomputed from eight class metrics.json files; noise_guard_summary.fired=0 was checked in every corresponding run_manifest.json.

\subsection{Few-shot 주 비교 (Few-shot Main Benchmark)}
\label{sec:experiments:main}

\begin{table*}[t]
\centering
\scriptsize
\setlength{\tabcolsep}{2.5pt}
\caption{MVTec AD 2 \texttt{test\_public} 해상도 설정값 1120 기준 few-shot 비교 (\%). 각 값은 3개 시드 평균 $\pm$ 표본 표준편차다. ScopeAD의 seed는 0, 1, 2이며 각 seed에서 8개 클래스를 평균한 뒤 시드 간 통계를 계산한다. $\dagger$ RegAD는 공식 체크포인트의 STN 제약으로 224 측정값을 병기했다. 굵은 글씨는 최고 평균값, 밑줄은 차순위다.}
\label{tab:mvtec_ad2_fewshot}
\begin{tabular}{r l r r r r r r r r}
\toprule
Shot & Metric & PatchCore & AnomalyDINO & FastRecon & RAD & SubspaceAD & SuperAD & RegAD (224)$^\dagger$ & Ours (ScopeAD) \\
\midrule
1 & AUPRO@0.05 & 38.45 $\pm$ 1.85 & 35.95 $\pm$ 2.24 & 14.92 $\pm$ 2.27 & 26.42 $\pm$ 1.71 & 44.87 $\pm$ 1.87 & \underline{46.31 $\pm$ 3.73} & 12.73 $\pm$ 0.38 & \textbf{51.31 $\pm$ 4.01} \\
   & pAUROC@0.05 & 64.19 $\pm$ 0.68 & \underline{70.67 $\pm$ 0.69} & 55.70 $\pm$ 0.67 & 59.57 $\pm$ 0.35 & 70.56 $\pm$ 0.56 & 70.13 $\pm$ 1.14 & 57.76 $\pm$ 0.30 & \textbf{72.93 $\pm$ 1.69} \\
   & I-AUROC & 65.51 $\pm$ 1.91 & 66.83 $\pm$ 1.08 & 58.10 $\pm$ 4.46 & 58.68 $\pm$ 3.03 & 67.02 $\pm$ 2.27 & \underline{67.68 $\pm$ 0.32} & 57.73 $\pm$ 1.82 & \textbf{73.47 $\pm$ 1.23} \\
\midrule
2 & AUPRO@0.05 & 42.25 $\pm$ 2.08 & 38.20 $\pm$ 1.90 & 16.73 $\pm$ 2.39 & 31.61 $\pm$ 1.47 & 45.90 $\pm$ 0.75 & \underline{49.80 $\pm$ 2.33} & 15.05 $\pm$ 0.30 & \textbf{56.03 $\pm$ 2.99} \\
   & pAUROC@0.05 & 65.43 $\pm$ 0.38 & \underline{72.26 $\pm$ 0.56} & 57.14 $\pm$ 1.03 & 61.78 $\pm$ 0.15 & 71.53 $\pm$ 0.18 & 71.52 $\pm$ 0.58 & 58.76 $\pm$ 0.49 & \textbf{75.86 $\pm$ 0.59} \\
   & I-AUROC & 67.53 $\pm$ 1.55 & 68.95 $\pm$ 3.84 & 57.63 $\pm$ 6.82 & 59.37 $\pm$ 2.54 & \underline{70.42 $\pm$ 1.79} & 70.19 $\pm$ 2.16 & 60.56 $\pm$ 3.94 & \textbf{74.93 $\pm$ 1.54} \\
\midrule
4 & AUPRO@0.05 & 46.94 $\pm$ 1.00 & 40.32 $\pm$ 0.75 & 21.88 $\pm$ 1.73 & 36.51 $\pm$ 1.02 & 47.14 $\pm$ 0.76 & \underline{54.25 $\pm$ 0.86} & 17.40 $\pm$ 1.07 & \textbf{58.73 $\pm$ 0.12} \\
   & pAUROC@0.05 & 66.90 $\pm$ 0.22 & 73.09 $\pm$ 0.91 & 59.26 $\pm$ 0.80 & 63.84 $\pm$ 0.41 & 72.45 $\pm$ 0.20 & \underline{73.59 $\pm$ 0.13} & 60.36 $\pm$ 0.48 & \textbf{76.94 $\pm$ 0.60} \\
   & I-AUROC & 68.62 $\pm$ 2.28 & 70.90 $\pm$ 0.97 & 59.73 $\pm$ 4.40 & 62.31 $\pm$ 4.84 & \underline{71.02 $\pm$ 0.56} & 70.97 $\pm$ 2.08 & 63.79 $\pm$ 2.34 & \textbf{75.88 $\pm$ 0.46} \\
\midrule
8 & AUPRO@0.05 & 49.98 $\pm$ 1.90 & 41.64 $\pm$ 1.01 & 27.32 $\pm$ 0.14 & 40.79 $\pm$ 0.79 & 47.28 $\pm$ 0.40 & \underline{55.44 $\pm$ 1.21} & 19.25 $\pm$ 0.64 & \textbf{59.74 $\pm$ 1.03} \\
   & pAUROC@0.05 & 68.22 $\pm$ 0.53 & 73.71 $\pm$ 0.20 & 61.74 $\pm$ 0.50 & 65.62 $\pm$ 0.69 & 72.56 $\pm$ 0.21 & \underline{74.31 $\pm$ 0.38} & 62.02 $\pm$ 0.32 & \textbf{77.51 $\pm$ 0.17} \\
   & I-AUROC & 69.20 $\pm$ 0.75 & \underline{71.29 $\pm$ 1.69} & 63.46 $\pm$ 4.29 & 65.84 $\pm$ 0.43 & 70.46 $\pm$ 1.18 & 70.75 $\pm$ 0.58 & 64.34 $\pm$ 1.05 & \textbf{75.99 $\pm$ 0.58} \\
\midrule
16 & AUPRO@0.05 & 52.73 $\pm$ 1.36 & 42.30 $\pm$ 1.25 & 35.23 $\pm$ 0.57 & 42.58 $\pm$ 1.10 & 47.63 $\pm$ 0.14 & \underline{55.89 $\pm$ 1.07} & 20.28 $\pm$ 0.80 & \textbf{60.06 $\pm$ 0.88} \\
   & pAUROC@0.05 & 69.36 $\pm$ 0.51 & 73.97 $\pm$ 0.23 & 65.63 $\pm$ 0.43 & 66.62 $\pm$ 0.89 & 72.78 $\pm$ 0.07 & \underline{74.60 $\pm$ 0.52} & 63.46 $\pm$ 0.76 & \textbf{77.61 $\pm$ 0.37} \\
   & I-AUROC & 69.00 $\pm$ 0.53 & \underline{72.07 $\pm$ 1.27} & 67.45 $\pm$ 1.16 & 67.60 $\pm$ 1.35 & 70.46 $\pm$ 1.07 & 69.79 $\pm$ 0.52 & 65.64 $\pm$ 1.34 & \textbf{75.92 $\pm$ 0.82} \\
\bottomrule
\end{tabular}
\end{table*}

표~\ref{tab:mvtec_ad2_fewshot}는 MVTec AD 2 \texttt{test\_public}에서의 few-shot 주 비교 결과다.
ScopeAD는 평가한 모든 shot에서 SuperAD보다 높은 평균 AUPRO를 보였다. 여기서 평균 차이는 관측된 효과이며, 통계적 유의성 검정의 결과는 아니다.
특히 16-shot에서 3-seed 평균 \textbf{60.06\,\% \AUPRO}를 달성하여 SuperAD(55.89\,\%) 대비 \textbf{+4.17\,pp}, pAUROC@0.05에서 \textbf{+3.01\,pp} 높은 성능을 기록했다.
1-shot AUPRO는 $51.31\pm4.01\,\%$로, seed 0의 $54.81\,\%$보다 평균이 낮고 서포트 선택에 따른 변동도 크다. 소수 참조에 대한 평균 성능과 표본 선택 민감도를 함께 고려해야 한다.

\begin{table*}[t]
\centering
\scriptsize
\setlength{\tabcolsep}{5pt}
\caption{MVTec AD 및 VisA few-shot 교차 데이터셋 비교 (Seed 0 단일 실행). 각 셀은 $\text{I-AUROC} \ / \ \text{P-AUROC}$ (전체 FPR 구간, \%) 순서다. 굵은 글씨는 최고값, 밑줄은 차순위다. ---는 해당 조건의 결과를 보고하지 않음을 뜻하며 순위에서 제외한다.}
\label{tab:mvtec_visa_fewshot}
\begin{tabular}{l r r r r r r r r r}
\toprule
Dataset & Shot & PatchCore & AnomalyDINO & FastRecon & RAD & SubspaceAD & SuperAD & RegAD & Ours (ScopeAD) \\
\midrule
MVTec AD & 1 & 88.12 / 96.22 & 92.52 / 94.49 & 88.84 / 96.60 & 80.83 / 93.25 & 89.25 / \underline{96.92} & \underline{92.54} / \textbf{96.93} & --- / --- & \textbf{92.76} / 96.77 \\
         & 2 & 92.67 / 96.97 & 95.17 / 95.56 & 93.51 / 97.29 & 84.15 / 94.27 & 90.95 / 97.24 & \underline{95.43} / \underline{97.36} & --- / --- & \textbf{96.13} / \textbf{97.39} \\
         & 4 & 95.06 / 97.37 & 95.92 / 96.07 & 95.61 / \textbf{97.87} & 85.38 / 94.67 & 91.35 / 97.16 & \underline{96.65} / \underline{97.70} & --- / --- & \textbf{96.99} / 97.68 \\
         & 8 & 96.13 / 97.80 & 96.67 / 96.45 & 96.76 / \textbf{98.16} & 87.85 / 95.29 & 92.94 / 97.17 & \underline{97.39} / \underline{97.95} & --- / --- & \textbf{97.60} / 97.93 \\
         & 16 & 98.15 / 97.99 & 97.48 / 96.80 & 98.52 / \textbf{98.38} & 90.29 / 95.61 & 92.51 / 97.19 & \textbf{98.94} / \underline{98.09} & --- / --- & \underline{98.74} / 98.05 \\
\midrule
VisA     & 1 & 86.09 / 95.56 & 87.06 / 95.62 & 84.65 / 96.75 & 72.67 / 95.37 & 85.44 / \textbf{98.09} & \textbf{90.29} / \underline{97.06} & 71.12 / 93.41 & \underline{89.88} / 96.73 \\
         & 2 & 87.87 / 96.29 & 90.64 / 96.63 & 87.67 / \underline{97.43} & 77.24 / 96.37 & --- / --- & \textbf{92.20} / \textbf{97.68} & 71.59 / 94.59 & \underline{92.15} / 97.32 \\
         & 4 & 89.52 / 96.92 & 91.37 / 97.57 & 89.21 / \underline{97.91} & --- / --- & --- / --- & \underline{92.94} / \textbf{98.02} & 76.42 / 96.14 & \textbf{94.19} / 97.82 \\
         & 8 & 92.36 / 97.18 & 93.11 / 97.99 & 91.53 / 98.11 & --- / --- & --- / --- & \underline{94.50} / \textbf{98.29} & 80.19 / 96.64 & \textbf{95.01} / \underline{98.25} \\
         & 16 & 93.64 / 97.30 & 93.53 / 98.30 & 93.38 / 98.32 & --- / --- & --- / --- & \underline{95.19} / \underline{98.43} & 81.22 / 96.98 & \textbf{95.99} / \textbf{98.49} \\
\bottomrule
\end{tabular}
\end{table*}

표~\ref{tab:mvtec_visa_fewshot}는 기존 MVTec AD 및 VisA 벤치마크에서의 교차 평가 결과다.
이미지 AUROC는 보고된 비교군 중 두 데이터셋의 모든 shot에서 최고 또는 차순위다.
반면 pixel AUROC는 MVTec AD 1/4/8/16-shot에서 3위이며 VisA 1-shot에서는 4위로, 지표별 우위가 일치하지 않는다.
표~\ref{tab:multiseed_visa_mvtec}는 2-shot과 16-shot의 세 서포트 선택에 따른 변동을 보여준다. 표의 0.00은 소수 둘째 자리까지 표시한 값이며 변동이 없음을 뜻하지 않는다.
이 비교는 전체 FPR 구간의 성능을 평가하므로 AD2의 low-FPR 결과와 구분한다.

\begin{table}[t]
\centering
\scriptsize
\setlength{\tabcolsep}{4pt}
\caption{ScopeAD 3-seed 교차 데이터셋 성능 (\%, mean $\pm$ sample SD).}
\label{tab:multiseed_visa_mvtec}
\begin{tabular}{l c c}
\toprule
Dataset & 2-shot (I-AUROC / P-AUROC) & 16-shot (I-AUROC / P-AUROC) \\
\midrule
MVTec AD & 96.04 $\pm$ 0.30 / 97.38 $\pm$ 0.04 & 98.52 $\pm$ 0.20 / 98.05 $\pm$ 0.00 \\
VisA     & 92.38 $\pm$ 1.04 / 97.40 $\pm$ 0.08 & 95.78 $\pm$ 0.18 / 98.50 $\pm$ 0.02 \\
\bottomrule
\end{tabular}
\end{table}

\subsection{관측 구성: 해상도와 매칭 방식}
\label{sec:experiments:observations}

\begin{table}[t]
\centering\scriptsize\setlength{\tabcolsep}{4pt}
\caption{단일 대 다중 해상도 비교 (linear0125, AD2, seed 0, \%).}
\label{tab:linear_res_ad2}
\begin{tabular}{l c c}
\toprule
해상도 구성 & 2-shot (AUPRO / pAUROC) & 16-shot (AUPRO / pAUROC) \\
\midrule
448 단독  & 51.14 / 75.80 & 52.01 / 76.52 \\
672 단독  & 54.38 / 74.97 & 55.01 / 76.17 \\
896 단독  & 52.65 / 73.46 & 55.12 / 74.62 \\
1120 단독 & 54.27 / 72.47 & 57.47 / 74.86 \\
\textbf{4개 해상도 결합} & \textbf{59.33 / 76.48} & \textbf{59.57 / 77.46} \\
\bottomrule
\end{tabular}
\end{table}

\noindent\textbf{동일 서포트의 다해상도 관측.}
표~\ref{tab:linear_res_ad2}에서 네 해상도 결합은 최선 단일 해상도보다 AUPRO가 2-shot에서 $+4.95$\,pp, 16-shot에서 $+2.10$\,pp 높다.
448 단독과의 16-shot 차이 $+7.56$\,pp는 고해상도 사용과 다해상도 결합을 함께 포함하므로, 다해상도 구성의 근거는 최선 단일 해상도와 용량 일치 대조군을 중심으로 해석한다.
이 비교는 정상 이미지 수를 늘리지 않고 입력의 공간적 관측 조건을 확장하는 선택을 지지한다. 추가 관측이 독립적인 정상 표본을 제공한다는 뜻은 아니다.

\noindent\textbf{특징 연결과 관측별 매칭.}
표~\ref{tab:component_ablation}의 layer concat 구성은 계층 특징을 연결한 뒤 하나의 최근접 이웃을 찾는다.
ScopeAD는 계층마다 이웃을 찾고 정상 참조 기반 비대칭 결합을 적용하며, 16-shot AUPRO가 concat보다 $2.09$\,pp 높다.
이 비교에서는 이웃 선택, 거리 기준화 및 계층 가중 방식이 함께 달라진다. 따라서 이는 전체 구성의 비교이며, bank 분리만의 독립 효과로 해석하지 않는다.

\subsection{용량 및 계산 비용}
\label{sec:experiments:cost}

\begin{table*}[t]
\centering\small\setlength{\tabcolsep}{5pt}
\caption{현재 선형 결합의 용량 및 비용 비교 (AD2 16-shot, seed 0, 8-class macro). 토큰 수는 이미지당·계층당·배그당 값이다. 지표는 AUPRO@0.05 / pAUROC@0.05 (\%). 비용은 클래스별 측정값의 평균이며, 시간은 피팅 초(s) 및 쿼리당 밀리초(ms), 메모리는 쿼리 단계 PyTorch peak allocated GiB다. 전체 GPU 메모리 사용량과는 다르다.}
\label{tab:resolution_control}
\begin{tabular}{l r c r r r}
\toprule
해상도 & 토큰 수 & AUPRO / pAUROC & Fit (s) & Query (ms) & Peak (GiB) \\
\midrule
4개 해상도 & 256 & 59.57 / 77.46 & 20.62 & 1526 & 3.07 \\
1120 단독 & 256 & 57.47 / 74.86 & 12.77 & 993 & 1.91 \\
1120 용량 일치 & 1024 & 57.14 / 75.49 & 14.42 & 1182 & 3.04 \\
4개 해상도 & 1024 & 59.17 / 77.56 & 27.04 & 1941 & 6.69 \\
\bottomrule
\end{tabular}
\end{table*}

표~\ref{tab:resolution_control}는 bank 용량과 실행 비용을 함께 비교한다.
4개 해상도에 각각 256개 토큰을 할당한 구성과 1120 단독에 1024개를 할당한 구성은 계층·배그당 총 bank 토큰 수가 같다.
다해상도 구성은 이 용량 일치 대조군보다 AUPRO가 $+2.44$\,pp 높지만, 쿼리 시간도 1182ms에서 1526ms로 증가한다.
따라서 bank 용량 통제와 계산량 통제를 구분해야 한다.
다해상도에서 토큰 수를 1024로 늘리면 AUPRO는 59.57에서 59.17로 감소하고 메모리 사용량은 증가하여, 더 큰 bank가 항상 유리하지는 않았다.

비용 측정은 FP16 autocast, batch size 1이며 CUDA 동기화를 포함한다.
각 클래스의 네 구성은 같은 GPU에서 측정했고 클래스별 측정값을 동일 가중 평균했다.
피팅 시간에는 서포트 디코딩, 특징 추출, bank 구축 및 척도 추정이 포함되며 모델 로딩, 서포트 선택 및 1회 warmup은 제외한다.
쿼리 시간에는 예측과 후처리가 포함된다. 이 값은 해당 실행 환경의 측정이며 비교 방법 전체의 속도 순위를 의미하지 않는다.

\subsection{정상 참조 기반 비대칭 결합의 대안 비교}
\label{sec:experiments:ablation}

\begin{table}[t]
\centering\scriptsize\setlength{\tabcolsep}{3pt}
\caption{결합 대안 및 관측 구성 비교 (AD2 \texttt{test\_public} 16-shot, seed 0, 8-class macro, \%). Full은 관측별 bank, q95, $\lambda=0.125$, 네 해상도를 사용한다. 결합 대안은 거리를 기준화하는 방식 또는 계층 가중치를 변경한다. 관측 구성 비교의 concat은 매칭과 결합을 함께 변경한다. $\Delta$는 Full 대비 AUPRO 차이(pp)이며 두 지표의 FPR 상한은 0.05다.}
\label{tab:component_ablation}
\begin{tabular}{l r r r}
\toprule
구성 & AUPRO & pAUROC & $\Delta$AUPRO \\
\midrule
\textbf{Full: q95, $\lambda=0.125$} & \textbf{59.57} & 77.46 & 0.00 \\
\multicolumn{4}{l}{\textit{결합 대안}} \\
정규화 제거 & 59.17 & \textbf{77.65} & $-0.40$ \\
계층 동일가중 ($\lambda=0.5$) & 58.21 & 76.31 & $-1.36$ \\
MAD 정규화 & 58.52 & 77.17 & $-1.05$ \\
\midrule
\multicolumn{4}{l}{\textit{관측 구성 비교}} \\
Layer concat bank & 57.48 & 74.80 & $-2.09$ \\
448 단일 해상도 & 52.01 & 76.52 & $-7.56$ \\
\bottomrule
\end{tabular}
\end{table}

식~\eqref{eq:linear-fusion}의 결합 계수 $w_{l,r,b}=\alpha_l/c_{l,r,b}$에 대해, 계층 가중치 $\alpha_l$과 정상 참조 척도 $c_{l,r,b}$를 각각 변경한 대안을 비교한다.
표~\ref{tab:component_ablation}은 대표 조건인 16-shot 결과를, 부록~\ref{sec:appendix:fusion}은 전체 가중치 및 shot별 정규화 결과를 제시한다.

\noindent\textbf{비대칭 가중과 동일가중.}
$\lambda=0.5$는 세 계층의 정규화 점수에 같은 계수를 부여하며, 전체 평균과 상수배 관계다.
$\lambda=0.125$는 이 동일가중 대안보다 AUPRO가 16-shot에서 $1.36$\,pp, 개발에 사용한 2-shot에서 $4.41$\,pp 높았다.
검사한 양의 가중치 중에서는 $0.125$가 두 shot에서 가장 높았으나, 16-shot에서 $0.25$와의 차이는 $0.09$\,pp로 작다(표~\ref{tab:lambda_sensitivity}).
이 결과는 깊은 계층에 더 큰 비중을 두는 결합을 지지한다. Deep-only인 $\lambda=0$ 비교가 없으므로 얕은 계층 추가 자체의 이득이나 정상 반응의 선택적 제거를 입증하지는 않는다.

\noindent\textbf{정상 참조 척도의 선택.}
q95는 MAD보다 모든 shot에서 AUPRO가 $0.76$--$1.69$\,pp 높았다(표~\ref{tab:norm_factorial_ad2}).
반면 미정규화 대비 AUPRO 차이는 1/2/4/8/16-shot 순서로 $-0.25$, $+0.27$, $-0.33$, $-0.01$, $+0.40$\,pp이며, pAUROC에서도 일관된 개선은 없다.
q95는 관측별 거리를 정상 참조 단위로 표현하여 비대칭 결합 계수를 정의하는 선택으로 사용한다. 이러한 구조적 역할과 식~\eqref{eq:scale-invariance}의 척도 불변성은 실제 성능 향상을 보장하지 않는다.

\subsection{획득 조건별 진단}
\label{sec:experiments:shift}

\begin{table}[t]
\centering\small
\caption{ScopeAD 선형 결합의 AD2 획득 조건별 성능 (16-shot, seed 0, \%). 네 조건이 모두 존재하는 동일한 7개 클래스(can, fabric, rice, sheet\_metal, vial, wallplugs, walnuts)의 macro이며 fruit\_jelly는 제외한다. 각 shift 번호는 데이터셋의 클래스별 획득 조건을 따른다.}
\label{tab:acquisition-strata}
\begin{tabular}{l r r}
\toprule
조건 & AUPRO@0.05 & pAUROC@0.05 \\
\midrule
Regular & 61.76 & 78.16 \\
Shift 1 & 60.48 & 78.10 \\
Shift 2 & 58.10 & 76.64 \\
Shift 3 & 59.45 & 77.04 \\
\bottomrule
\end{tabular}
\end{table}

표~\ref{tab:acquisition-strata}는 같은 클래스 집합에서 획득 조건별 성능을 분리한 진단이다.
Regular AUPRO 61.76에 비해 Shift 1--3은 58.10--60.48로 낮아, 고정 결합 이후에도 조건 변화에 따른 저하가 남는다.
이 결과는 공개 개발 세트 내부의 분석이며 미관측 환경에 대한 독립 검증이나 비교 방법 대비 shift 우위를 뜻하지 않는다.
% Legacy gated private results and fig_acquisition_shift_by_shot.pdf are excluded: they do not evaluate the current linear method.

\subsection{합성 열화 조건 비교}
\label{sec:experiments:corruption}

\begin{table*}[t]
\centering\scriptsize\setlength{\tabcolsep}{4pt}
\caption{MVTec AD 합성 열화 비교 (15개 클래스, 16-shot, seed 0, 열화 강도 3, 서포트는 열화 없음). 각 셀은 AUPRO@0.05 / pAUROC@0.05 (\%)이며, 굵은 글씨는 각 지표의 최고값, 밑줄은 차순위다. $\dagger$ FastRecon은 DINOv2-L/14 기반 국소 이웃 능형 복원 변형이며 원 논문의 재현 결과가 아니다. $\ddagger$ RegAD는 공식 구현과 구별되는 DINOv2 밀집 매칭 변형이다. RAD는 로컬 DINOv3 ViT-B/16 캐시 경로로 실행했다.}
\label{tab:corruption}
\begin{tabular}{l c c c c c}
\toprule
방법 & Clean & Brightness & Contrast & Defocus blur & Gaussian noise \\
\midrule
PatchCore & 83.97 / 88.80 & 83.43 / \underline{88.38} & 77.15 / 85.27 & 81.58 / 87.41 & \textbf{74.03 / 83.66} \\
AnomalyDINO & 77.62 / 85.31 & 72.54 / 82.61 & 72.08 / 82.48 & 70.69 / 81.06 & 33.02 / 66.81 \\
FastRecon$^\dagger$ & \textbf{85.74 / 90.82} & 77.53 / 86.35 & 69.72 / 84.47 & \textbf{82.35 / 88.90} & 57.70 / 77.37 \\
RAD & 80.12 / 87.03 & 78.70 / 86.04 & 74.96 / 84.33 & 70.44 / 81.13 & 49.23 / 72.05 \\
SubspaceAD & 83.89 / 86.98 & 80.45 / 84.91 & 78.49 / 83.39 & 78.17 / 83.39 & 64.61 / 78.01 \\
SuperAD & 84.49 / \underline{89.15} & \underline{83.68} / \textbf{88.47} & \underline{82.27} / \textbf{87.86} & \underline{82.17} / \underline{87.74} & \underline{71.10} / \underline{82.45} \\
RegAD$^\ddagger$ & 38.11 / 78.65 & 36.14 / 77.02 & 32.01 / 74.77 & 37.33 / 77.91 & 28.12 / 70.00 \\
\midrule
Ours (ScopeAD) & \underline{85.42} / 88.98 & \textbf{84.64} / 88.35 & \textbf{83.40} / \underline{87.59} & 81.46 / 86.38 & 68.43 / 80.19 \\
\bottomrule
\end{tabular}
\end{table*}

표~\ref{tab:corruption}은 Hendrycks와 Dietterich~\cite{hendrycks2019benchmarking}의 합성 열화 평가 방식을 따라 테스트 이미지에만 열화를 적용한 결과다.
ScopeAD의 AUPRO는 Brightness와 Contrast에서 비교군 중 가장 높고, Clean에서는 FastRecon 다음이다. 반면 Defocus blur에서는 FastRecon과 SuperAD보다, Gaussian noise에서는 PatchCore와 SuperAD보다 낮다.
Gaussian noise에서 ScopeAD의 AUPRO는 Clean 대비 16.99\,pp 낮아져, PatchCore의 9.94\,pp와 SuperAD의 13.39\,pp보다 하락 폭이 크다.
pAUROC에서는 어느 조건에서도 최고값이 아니며, Contrast에서만 차순위다.
따라서 이 결과는 ScopeAD가 합성 열화 전반에서 우위라는 근거가 아니라, 밝기·대비 변화에서의 강점과 잡음·초점 열화에서 남는 취약성을 함께 보여준다.
서포트에 열화를 적용하지 않은 단일 seed 평가이므로 조건 간 차이의 통계적 유의성은 검정하지 않았다. 비교군의 해상도와 백본 설정은 방법별로 다르다.

\subsection{공식 private 평가}
\label{sec:experiments:private}

\begin{table}[t]
\centering\scriptsize\setlength{\tabcolsep}{3pt}
\caption{MVTec AD 2 공식 private 평가 (8-class macro, \%). 각 셀은 \texttt{test\_private} / \texttt{test\_private,mix} 순서다. ScopeAD는 16-shot, seed 0 단일 제출이다. ClassF1과 SegF1의 임계값은 서포트를 제외한 \texttt{train/good}에서 정한다. SuperAD 행은 전체 학습 데이터 기준 게시값이며 few-shot 결과와 직접 순위화하지 않는다.}
\label{tab:private}
\begin{tabular}{l c c c}
\toprule
방법 & AUPRO@0.05 & ClassF1 & SegF1 \\
\midrule
ScopeAD (16-shot) & 65.21 / 63.25 & 76.57 / 77.68 & 37.71 / 34.00 \\
\midrule
SuperAD (전체 데이터)~\cite{superad} & 60.51 / 58.37 & 70.20 / 74.40 & 47.18 / 42.51 \\
\bottomrule
\end{tabular}
\end{table}

표~\ref{tab:private}은 공식 서버에 제출한 ScopeAD의 private 결과다(제출 상태 DONE, 2026-10-14 확인).
임계값은 서포트 제외 \texttt{train/good}의 8번째 간격 표본에서 클래스별 정상 픽셀 p95의 1.421배로 정했으며 private 레이블은 사용하지 않았다.
AUPRO는 두 조명 조건에서 게시된 전체 데이터 SuperAD보다 높지만, SegF1은 낮다. 두 방법은 정상 데이터 예산, 백본, 입력 크기, 후처리가 달라 직접 비교가 아니다.
이 결과는 단일 seed의 한 번 제출이며 같은 조건의 few-shot 비교군 private 결과는 없다.
